\section{Introduction}\label{sec:introduction}

An algebraic extension of a number field is \emph{solvably closed} if it has no nontrivial abelian algebraic extension. Algebraic closures and maximal prosolvable extensions are examples. The isomorphism theorems of Neukirch and Uchida show that Galois groups of such fields retain remarkably precise arithmetic information. The corresponding question for open homomorphisms asks whether this reconstruction also recovers field embeddings.

Let $E_i/F_i$ be Galois extensions of number fields, with each $E_i$ solvably closed, and write $G_i=\Gal(E_i/F_i)$. A field embedding $j:E_2\hookrightarrow E_1$ \emph{induces} a homomorphism $\alpha:G_1\to G_2$ if
\[
 g\circ j=j\circ\alpha(g)\qquad(g\in G_1).
\]
In particular, such an embedding sends $F_2$ into $F_1$. Uchida conjectured that every continuous homomorphism with open image arises in this way \cite[p.~595]{uchida1981}. We prove the conjecture in its full form.

\begin{theorem}\label{thm:main}
Let $F_1,F_2$ be number fields, and let $E_i/F_i$ be possibly infinite solvably closed Galois extensions. Every continuous open homomorphism
\[
 \alpha:\Gal(E_1/F_1)\longrightarrow\Gal(E_2/F_2)
\]
is induced by a unique field embedding $j:E_2\hookrightarrow E_1$.
\end{theorem}

The kernel is unrestricted. The proof establishes the necessary arithmetic compatibility from openness itself.

\subsection{Historical context}
The reconstruction of number fields from their Galois groups begins with Neukirch's work on normal number fields and decomposition groups \cite{neukirch1969a,neukirch1969b}. The passage to arbitrary number fields and prescribed group isomorphisms involved Uchida, Iwasawa, and Ikeda; their contributions are recorded in the contemporary accounts \cite{uchida1976,neukirch1977}. Uchida subsequently established the isomorphism theorem for arbitrary solvably closed Galois extensions of number fields \cite{uchida1979}.

In the paper formulating the homomorphism conjecture, Uchida proved it when the source number field is $\Q$, proved uniqueness in general, and established existence under a suitable local condition on decomposition groups \cite[Theorem~1, Proposition~2, and Theorem~2]{uchida1981}. His local analysis also associates primes to one another whenever a chosen pro-$\ell$ inertia group survives under the homomorphism, for an odd auxiliary prime $\ell$. For each fixed auxiliary prime $\ell$, this partial correspondence reaches all but finitely many target primes \cite[Proposition~1]{uchida1981}.

A later reduction of Hoshi identifies an exact arithmetic condition for existence: an open homomorphism is induced by a field embedding if and only if it preserves the full cyclotomic character \cite[Theorem~3.4]{hoshi2025}. Hoshi's subsequent work reformulates the unrestricted conjecture in terms of radicals and normality of the field fixed by a kernel \cite[Theorem~B]{hoshi2026}. Related work on finite solvable quotients includes the isomorphism theorem of Sa\"idi and Tamagawa \cite{saiditamagawa2022} and the cyclotomic-compatible homomorphism theorem of Mao and Sa\"idi \cite{maosaidi2025}. Our proof uses the cyclotomic criterion at its endpoint. The main task is to derive that criterion from the partial local correspondences, even though their domains and exceptional sets initially depend on $\ell$.

\subsection{The main ideas}
Two arguments used below are useful beyond the final reduction. First, an additive $\ell$-adic character of the absolute Galois group of a number field is determined by its values on any infinite set of primes of absolute residue degree one. More precisely, a character vanishing at all those primes is zero. The logarithmic-rank input is a special case of the $\ell$-adic Waldschmidt--Masser theorem \cite{waldschmidt1981,masser1981,dasgupta2023}. We give a complete proof using Laurent's interpolation-determinant method \cite{laurent1991} and a translate-and-degree zero estimate on a torus. Its variants for Frobenius averages and corestriction permit us to treat primes of arbitrary residue degree as well. These statements concern the actual unit-logarithm image and require no hypothesis on its rank.

Second, we refine the Kummer argument in Uchida's proof of cofinite target coverage \cite[Proposition~1]{uchida1981}. After logarithmic cyclotomic compatibility has been established, the number of target primes missed by the partial correspondence can be bounded independently of $\ell$. A missed prime supplies a Kummer class. On killing a finite cyclotomic twist, these classes lie in one character component of an $S$-unit representation. The full extension degree may grow with $\ell$, but the multiplicity of that character is bounded by the number of orbits of $S$, hence by the degree of the fixed source field. This uniformity allows us to choose $\ell$ by simultaneous Kummer conditions and force enough corresponding primes to have the same rational residue characteristic.

Here is the route from these ingredients to Theorem~\ref{thm:main}. Put $P=\Gal(\Q^s/\Q)$, where $\Q^s$ is the maximal prosolvable extension of $\Q$. For a number field $F$, let $G_F$ denote its absolute Galois group. It suffices to identify every continuous open homomorphism $G_F\to P$ with ordinary restriction up to conjugation. Indeed, every solvably closed field contains $\Q^s$, and this identification will give the full cyclotomic compatibility required by Hoshi's theorem.

We first restrict to a normal open subgroup of $G_{\Q}$ and extend the map to the normalizer of its kernel. Write this extension as $\beta:G_k\to P$. Following the regular-module constructions used by Uchida \cite{uchida1979}, we enlarge finite quotients of $P$ by elementary abelian kernels consisting of many copies of a regular representation. Each quotient $R=\Gal(L/\Q)$ acts on the additive character space of $G_L$. Suitable characters distinguish a cyclic subgroup generated by a target Frobenius element from each of its proper subgroups. Comparing Frobenius averages proves that $\beta$ is surjective and preserves logarithmic cyclotomic characters for an infinite congruence class of auxiliary primes.

Next, pass to the preimage of $P_{\Q(i)}=\Gal(\Q^s/\Q(i))$. A sign character shows that, for almost every split target prime, only finitely many source primes can correspond to it as $\ell$ varies. The uniform Kummer bound then forces a correspondence in the same residue characteristic for almost every such target prime. Finally, use rational primes split completely in suitable finite extensions. Their additive Frobenius evaluations identify the two character pullbacks, one from $\beta$ and one from ordinary restriction. A faithful representation in the character space identifies the group maps modulo every finite quotient; compactness supplies a single conjugation.

Section~\ref{sec:preliminaries} states the anabelian inputs and proves the local norm identity. Section~\ref{sec:characters} develops the character tests. Section~\ref{sec:generation} obtains surjectivity and logarithmic cyclotomic compatibility. Section~\ref{sec:matching} proves the uniform bound and matches residue characteristics. Section~\ref{sec:identification} identifies the map to $P$ and proves Theorem~\ref{thm:main}.
